
\subsection{Homotopiegruppen}

\begin{definition}
  Sei \(A\) punktierter Typ und \(n:\N\).
  Dann ist
  \[\pi_n\colonequiv \|\Omega^n(A)\|_0\]
  --- was wir für \(n\geq 1\) \begriff{\(n\)-te Homotopiegruppe}\index{\(\pi_n(A)\)} von \(A\) nennen.
  Für eine punktierte Abbildung \(f:A\to_\ast B\), also \(f:A\to B\) mit \({p:f(\ast)=\ast}\)
  ist \(\pi_n(f)\)\index{\(\pi_n(f)\)} gegeben durch
  \begin{align*}
    \pi_n(f):\pi_n(A)&\to \pi_n(B) \\
    (l:\ast=_A\ast)&\mapsto p^{-1}\kon l\kon p
  \end{align*}
\end{definition}

Um den Namen zu rechtfertigen, sollten wir zeigen, dass es sich bei den \(\pi_n(A)\) um Gruppen handelt.
Tatsächlich ist \(\pi_n(A)\) für \(n\geq 2\) eine abelsche Gruppe, was wir zunächst als stärkeres Resultat für doppelte Schleifenräume zeigen und als Eckmann-Hilton Argument bekannt ist:

\begin{lemma}
  \label{eckmann-hilton}
  Sei \(A\) punktiert. Für \(x,y:\Omega^2(A)\) haben wir \(x\kon y=y\kon x\).
\end{lemma}

\begin{beweis}
  
\end{beweis}

\begin{bemerkung}
  Sei \(A\) punktierter Typ.
  \begin{enumerate}[(a)]
  \item Für \(n\geq 1\) ist \(\pi_n(A)\) eine Gruppe mit der durch Konkatenation von Gleichheiten gegebenen Verknüpfung.
  \item Für \(n\geq 2\) ist \(\pi_n(A)\) eine abelsche Gruppe\footnote{D.h.\ \(x\kon y=y\kon x\) für alle \(x,y\).}.
  \item Sei \(f:A\to_\ast B\). Für \(n\geq 1\) ist \(\pi_n(f)\) ein Gruppenhomomorphismus.
  \end{enumerate}
\end{bemerkung}

\begin{beweis}
  \begin{enumerate}[(a)]
  \item Wir haben die entsprechenden Gesetze bereits für allgemeine Identitätstypen \(x=y\) gezeigt.
    Nun ist \(\|\Omega^n(A)\|_0\) eine Menge für \(x,y,z:|\Omega^n(A)\|_0\) gilt zum Beispiel Assoziativität per \(\|\_\|_0\)-Induktion und Andwenden von \(|\_|_0\) auf die entsprechende Gleichheit in \(\Omega^n(A)\).
  \item \Cref{eckmann-hilton}.
  \item \(f(s\kon t)=p^{-1}\kon s\kon t\kon p=p^{-1}\kon s\kon p\kon p^{-1}\kon t\kon p=f(s)\kon f(t)\)
  \end{enumerate}
\end{beweis}
